\chapter{إعادة فحص الأقفال والتزامن المتقدم}
\label{CH:LOCK2}

بعد دراسة جميع مكونات النواة (الذاكرة، المصايد، الجدولة، والمقاطعات، ونظام الملفات)، يقدم هذا الفصل مراجعة شاملة وإعادة تفحص متعمقة لكيفية تفاعل الأقفال وميكانيكيات التزامن المتقدمة لحماية هياكل البيانات في \texttt{xv6}.

يقدم هذا الفصل تحليل تفصيلي لاستخدام الأقفال في المخزن المؤقت للكتل، والعقد الفهرسية، الأنابيب، والعمليات.

\section{أنماط الأقفال والتنسيق}
\label{sec:locking_patterns}

يحتوي المخزن المؤقت للكتل (Buffer Cache) في \texttt{kernel/bio.c} على مستويين من الأقفال:
\begin{enumerate}
  \item \textbf{القفل العام (\texttt{bcache.lock})}: يحمي القائمة المتصلة المزدوجة للمخزن المؤقت (\texttt{bcache.head}). يُستخدم فقط عند البحث عن كتلة أو إعادة استخدام كتلة حرة (LRU).
  \item \textbf{القفل الفردي للكتلة (\texttt{b->lock})}: قفل نوم (\texttt{sleeplock}) مخصص لكل كتلة مؤقتة حية. يحمي محتويات البيانات داخل الكتلة (\texttt{b->data}) ويمنع العمليات الأخرى من تعديل نفس الكتلة أثناء تنفيذ عمليات الإدخال والإخراج المادية على القرص.
\end{enumerate}

يسمح هذا التصميم للأنوية المختلفة بالقراءة والكتابة في كتل قرصية مختلفة بالتوازي دون حيازة القفل العام للمخزن المؤقت، مما يحسن الأداء بصورة ملحوظة.

\section{أنماط التزامن الشبيهة بالأقفال}
\label{sec:lock_like_patterns}

تتبع العقد الفهرسية (Inodes) نمطًا مشابهًا للمخزن المؤقت:
\begin{itemize}
  \item \textbf{القفل العام جدول العقد (\texttt{icache.lock})}: يحمي القائمة الموزعة للجدول وتخصيص المداخل في الذاكرة.
  \item \textbf{القفل الفردي للعقدة (\texttt{ip->lock})}: قفل نوم يحمي حقول العقدة الفهرسية ومحتويات كتلها المربوطة.
\end{itemize}

\section{كود بدون أقفال والتعليمات الذرية}
\label{sec:no_locks_at_all}

في بعض الحالات النادرة مثل السجل العتادي \texttt{tp} الخاص بمعرف وحدة المعالجة أو استخدام التعليمات الذرية في الجدولة، يمكن الاستغناء عن الأقفال الصريحة بشريط تعيين ذرية محكومة عتاديًا.

\section{التوازية والتوسع الأفقية}
\label{sec:parallelism}

حماية مصفوفة العمليات (\texttt{proc}) من أصعب وظائف التزامن داخل النواة لأن العملية يمكن أن تُعدل حالتها بواسطة عمليات أخرى (مثل \texttt{wait()}، \texttt{kill()}، و \texttt{exit()}).

تستخدم النواة القفل المخصص لكل عملية (\texttt{p->lock}):
\begin{itemize}
  \item يحمي حالة العملية (\texttt{p->state})، ومستقبل القناة (\texttt{p->chan})، وراية الموت (\texttt{p->killed})، ومعرف العملية (\texttt{p->pid}).
  \item يُحاز القفل عند القراءة أو التعديل على حالة العملية من قبل المجدول أو العمليات الأب.
\end{itemize}

\section{تمارين}
\label{sec:exercises_lock2}

\begin{enumerate}
  \item اشرح لماذا يُفضل استخدام أقفال النوم (Sleep-locks) لحماية كتل القرص بدلاً من أقفال التناوب (Spinlocks)؟
  \item تتبع ترتيب الأقفال المحوزة أثناء تنفيذ استدعاء النظام \texttt{sys\_link()} لمنع حدوث الجمود.
\end{enumerate}
